\section*{Capítulo I, Problema X}

\textbf{Problema:}

10. Sean $A, B$ dos vectores no nulos en el espacio $n$-dimensional. Sea $\theta$ el ángulo entre ellos. Si $\cos \theta = 1$, demuestra que $A$ y $B$ tienen la misma dirección. Si $\cos \theta = -1$, demuestra que $A$ y $B$ tienen direcciones opuestas. [Pista: Si $c$ es la componente de $A$ a lo largo de $B$, demuestra que $(A - cB)^2 = 0$].

\textbf{Solución:}

\begin{itemize}
    \item \textbf{Caso 1: $\cos \theta = 1$}

    Si $\cos \theta = 1$, entonces el ángulo entre $A$ y $B$ es $0$ grados. Esto implica que $A$ y $B$ están alineados en la misma dirección. Matemáticamente, esto significa que existe un escalar positivo $k > 0$ tal que $A = kB$.

    \item \textbf{Caso 2: $\cos \theta = -1$}

    Si $\cos \theta = -1$, entonces el ángulo entre $A$ y $B$ es $180$ grados. Esto implica que $A$ y $B$ están alineados en direcciones opuestas. Matemáticamente, esto significa que existe un escalar negativo $k < 0$ tal que $A = kB$.

    \item \textbf{Demostración de la pista:}

    Si $c$ es la componente de $A$ a lo largo de $B$, entonces $c = \frac{A \cdot B}{\|B\|^2}$. Por lo tanto:
    \[
        A - cB = A - \frac{A \cdot B}{\|B\|^2}B.
    \]

    Calculamos el cuadrado de la norma de $A - cB$:
    \begin{align*}
        \|A - cB\|^2 &= (A - cB) \cdot (A - cB) \\
        &= A \cdot A - 2c(A \cdot B) + c^2(B \cdot B).
    \end{align*}

    Sustituyendo $c = \frac{A \cdot B}{\|B\|^2}$, obtenemos:
    \begin{align*}
        \|A - cB\|^2 &= \|A\|^2 - 2\frac{(A \cdot B)^2}{\|B\|^2} + \frac{(A \cdot B)^2}{\|B\|^2} \\
        &= \|A\|^2 - \frac{(A \cdot B)^2}{\|B\|^2}.
    \end{align*}

    Si $\cos \theta = \pm 1$, entonces $A \cdot B = \pm \|A\|\|B\|$, lo que implica que $\|A - cB\|^2 = 0$. Por lo tanto, $A - cB = 0$, lo que demuestra que $A$ y $B$ están alineados.
\end{itemize}